\section{Finite-memory approximation uniform over bounded-cost policies}
\label{sec:r75-stopped-approximation}

Exact finite memory and approximation of an experiment at bounded
expected cost have different quantifiers.  The noisy construction gives
an explicit approximation of the undamped conditional process, using one
finite model for an entire class of adaptive stopping policies.  The
following statement is an upper bound; it does not claim an optimal
approximation exponent among arbitrary classical models.

\begin{theorem}[Stopped-law approximation from conditional damping]
\label{thm:r75-stopped-approximation}
Fix a finite angle command alphabet, initial qubit $I/2$, and the
undamped instruments \eqref{eq:r75-noisy-instrument} with $\alpha=1$.
Let $\mathcal P_B$ be any class of common randomized policies for the
next command, stopping and a terminal qubit POVM, such that the number
$\tau$ of physical calls satisfies
$\mathbb E_{1}^{\pi}\tau\le B<\infty$ for every $\pi\in\mathcal P_B$.
The policy has access to the visible record and its own classical
randomness; there is no additional quantum register entangled with the
qubit.

For every $0<\epsilon<1$, replace the instrument by the one with
$\alpha=1-\epsilon$ and use its common finite realization from
Theorem~\ref{thm:r75-noisy-memory}.  This one model, independently of the
chosen policy, obeys
\begin{equation}\label{eq:r75-stopped-tv}
 \sup_{\pi\in\mathcal P_B}
 \left\|\mathcal L_1^\pi-\mathcal L_{1-\epsilon}^\pi\right\|_{\rm TV}
 \le \epsilon B/8.
\end{equation}
Here the complete stopped record includes the stopping time, commands,
reports, terminal measurement result and a nontermination symbol.  The
stronger bound holds for the block-diagonal classical--quantum terminal
states before applying the terminal POVM; its nontermination block is
one-dimensional and records only the nontermination mass.
The simulator uses at most
$K(1-\epsilon)$ labels with a free mixed start, or
$K(1-\epsilon)+1$ with a counted pure start.

Thus for $0<\delta<B/8$, one may take $\epsilon=8\delta/B$ and use
\begin{equation}\label{eq:r75-approx-count}
 1+\max\left\{3,
       \left\lceil\frac\pi{\arccos(1-8\delta/B)}\right\rceil\right\}
   =\frac\pi4\sqrt{B/\delta}+O(1)
\end{equation}
pure-start labels as $\delta/B\downarrow0$.  For $\delta\ge B/8$ the
$\alpha=0$ model uses three pure-start labels and satisfies the same
error requirement.  Policy memory is separate: a policy with $L$
counted labels and a prescribed initial label can be composed with the
mixed core using at most $LK(1-\epsilon)+1$ labels, with its first actual
call processed during entry into the product core.
\end{theorem}
\begin{proof}
For every positive qubit operator $X$ and every angle,
\[
 \|\mathcal D(X)-U_\theta XU_\theta^*\|_1\le\tr X.
\]
Indeed, after division by $\tr X>0$, the difference from $I/2$ has
eigenvalues $\pm |r|/2$ for a Bloch vector of length at most one; the
zero-trace case is immediate.  The perturbation of branch zero is
$\epsilon[\mathcal D-\operatorname{Ad}_{U_\theta}]/8$, that of branch one
is zero, and that of branch two is the negative of branch zero.  Hence
the one-call flagged output satisfies
\begin{equation}\label{eq:r75-instrument-distance}
 \frac12\sum_{s=0}^2
   \|[\Psi_s^{1-\epsilon,\theta}-\Psi_s^{1,\theta}](X)\|_1
 \le\frac\epsilon8\tr X.
\end{equation}
This bound is uniform in the command.  Averaging it over a randomized
command choice preserves it.

First cap the experiment at a deterministic horizon $H$.  Retain the
classical history and the subnormalized qubit state in every live or
stopped branch.  A subsequent common policy step, including a terminal
measurement, is a positive trace-preserving map on this direct sum, and
therefore contracts the trace norm of Hermitian differences.  To see
this directly, write such a difference as $Z_+-Z_-$, with
$Z_\pm\ge0$.  Positivity, trace preservation and the triangle inequality
give $\|\Phi(Z)\|_1\le\tr Z_++\tr Z_-=\|Z\|_1$.

Telescope between the undamped and damped $H$-step experiments, replacing
one physical call at a time.  In the term for call $t+1$, keep the first
$t$ calls undamped and the remaining tail damped.  The input to the
changed call is then the target positive live state of total trace
$\mathbb P_1^\pi\{\tau>t\}$.  Apply
\eqref{eq:r75-instrument-distance} there and contract the difference
through its common tail.  The resulting terminal trace distance is at
most
\begin{equation}\label{eq:r75-capped-distance}
 \frac\epsilon8\sum_{t=0}^{H-1}
      \mathbb P_1^\pi\{\tau>t\}
 =\frac\epsilon8\mathbb E_1^\pi\min\{\tau,H\}.
\end{equation}
The reference distribution in this bound is the undamped target, not an
uncontrolled distribution of a hybrid or approximate process.

To pass to the uncapped experiment, first retain separately every stopped
history of length at most $H$ and merge the unfinished branches into one
symbol.  Coarse graining cannot increase the distance.  As $H$ increases,
the sum of the trace norms of the differences on the finite stopped
blocks increases to its full sum, while the unfinished probabilities
decrease to the respective nontermination probabilities.  Thus the
left side converges to the stated stopped classical--quantum distance,
including the absolute difference at the nontermination symbol.  The
right side increases to $\epsilon\mathbb E_1^\pi\tau/8$.  A terminal
POVM and further classical processing contract this distance, proving
\eqref{eq:r75-stopped-tv}.  The argument permits all the policy
randomization because that randomness is simply part of the common
classical transition at each step.

The noisy dictionary reproduces these complete subnormalized word states
exactly, so replacing the noisy quantum system by its classical
realization introduces no further error.  The count follows from the
preceding theorem.  Substitution of $\epsilon=8\delta/B$ proves
\eqref{eq:r75-approx-count}; use
$\arccos(1-x)=\sqrt{2x}+O(x^{3/2})$ as $x\downarrow0$.
The same perturbation bound is valid at $\epsilon=1$, where the exact
two-state core was given explicitly.  Finally a product with the
$L$-label policy has $LK$ core labels.  Its separate initial label has
the initial terminal decoder and, on the first call, the appropriate
policy-weighted transition into that product core.  This proves the
claimed count without omitting the first physical call.
\end{proof}

The bound is uniform over a class defined by the target's expected number
of calls.  It does not imply that the approximate experiment has expected
cost at most $B$: total variation controls bounded losses, whereas an
unbounded stopping time requires a separate uniform integrability or
tail hypothesis.  In particular the possible nontermination mass of the
approximate process is bounded by $\epsilon B/8$ but need not vanish.
If a deterministic horizon is part of the policy class, both processes
terminate within that horizon and \eqref{eq:r75-capped-distance} applies
directly.  No extra quantum register, quantum input encoding, precision
budget, or external policy memory is included in the stated simulator
label count.
If $B=0$, the target makes no call almost surely.  The same policy then
also makes no call in the simulator, and a single label with the initial
terminal decoder suffices with zero error.

\subsection{Status of the inherited arithmetic consequences}
\label{sec:r75-arithmetic-lift}

For completeness, the claim after Theorem~\ref{thm:r75-noisy-memory}
uses only the two implications already proved in
Lemma~\ref{lem:r75-full-instrument-bounds}.  With
$\alpha=1-\epsilon$, $x=\theta/(2\pi)$ and
\[
 \mathfrak A_x(\epsilon)
   =\min_{q\ge1}\max\{3,q,\|qx\|/(q\epsilon)\},
\]
the lower proof of Theorem~\ref{thm:r71-envelope} applies because
$\alpha e^{i\theta}\in\mathcal K_M$ is necessary for a full-instrument
$M$-state realization.  Its upper proof constructs a regular polygon
containing $\alpha e^{i\theta}$, so it also constructs the complete
instrument, rather than just its scalar decoder.  Therefore
\[
 \tfrac12\mathfrak A_x(\epsilon)
 \le N_{\mathrm{mix}}(1-\epsilon,\theta)
 \le (1+4\pi^2)\mathfrak A_x(\epsilon).
\]
The exact $+1$ pure-start increment does not change logarithmic
exponents, divergence, or the constant-factor cubic-root criterion.
All subsequent number-theoretic reasoning is inherited from
Section~\ref{sec:r71-arithmetic}; no new irrationality exponent is
assigned to a particular algebraic rotation.  In particular the full
Karpelevi\v c region alone has not been shown sufficient for this complete
instrument away from the sectors treated above.
